% Created 2026-06-27 Sat 14:10
% Intended LaTeX compiler: pdflatex
\documentclass[11pt,letterpaper]{article}
\usepackage[utf8]{inputenc}
\usepackage[T1]{fontenc}
\usepackage{graphicx}
\usepackage{longtable}
\usepackage{wrapfig}
\usepackage{rotating}
\usepackage[normalem]{ulem}
\usepackage{amsmath}
\usepackage{amssymb}
\usepackage{capt-of}
\usepackage{hyperref}
\usepackage{amsmath}
\usepackage{amssymb}
\usepackage{geometry}
\geometry{left=1.2in, right=1.2in, top=1.0in, bottom=1.0in}
\usepackage{tcolorbox}
\tcbuselibrary{skins,breakable}
\newenvironment{defbox}{\begin{tcolorbox}[colback=green!3!white,colframe=green!45!black,title={\small\textit{Foundational definition}},breakable]}{\end{tcolorbox}}
\usepackage{fancyhdr}
\usepackage{booktabs}
\usepackage{array}
\usepackage{parskip}
\usepackage{microtype}
\PassOptionsToPackage{hidelinks}{hyperref}
\setlength{\parindent}{0pt}
\setlength{\parskip}{6pt}
\fancypagestyle{plain}{\fancyhf{}\fancyfoot[C]{\thepage}\renewcommand{\headrulewidth}{0pt}}
\fancypagestyle{body}{%
\fancyhf{}%
\fancyhead[C]{\small Document 1 of 5 $\cdot$ DOI: 10.5281/zenodo.18927154 $\cdot$ CC BY 4.0}%
\fancyfoot[C]{\thepage}%
\renewcommand{\headrulewidth}{0.4pt}}
\author{Bradley K. DePuy, Independent Researcher}
\date{2026}
\title{The Medium: Foundation Definition}
\hypersetup{
 pdfauthor={Bradley K. DePuy, Independent Researcher},
 pdftitle={The Medium: Foundation Definition},
 pdfkeywords={},
 pdfsubject={},
 pdfcreator={Emacs 29.3 (Org mode 9.6.15)}, 
 pdflang={English}}
\begin{document}

\pagestyle{plain}
\begin{titlepage}
\vspace*{4cm}
\begin{center}
  {\LARGE\bfseries The Medium}\\[0.6em]
  {\large An Attempted Description of a Pre-Geometric Substrate --- Foundation Definition}\\[2.5em]
  {\normalsize Bradley K. DePuy, Independent Researcher}\\[0.5em]
  {\normalsize ORCID: \mbox{\texttt{0009-0001-0328-2299}}}\\[0.5em]
  {\normalsize 2026}
\end{center}
\end{titlepage}
\pagestyle{body}

\section{Definition}
\label{sec:org770c219}

\subsection{The Substrate}
\label{sec:orgd3d4d6d}

The Medium is the substrate --- denoted \(M\). It has no beginning and no end, no geometry, no time, no parts, and no structure prior to distinction. \(M\) is not a space or a set of points. It is what is.

\subsection{The Primitive: \(\iota\)}
\label{sec:orga5d36b5}

The sole primitive is \(\iota\): a distinction within \(M\). Not an event --- event implies sequence, which implies time. \(\iota\) is prior to time. A distinction is the substrate differentiating within itself.

Distinction is self-grounding. Any attempt to explain what causes it already uses it. Causation, before, after --- these are all distinctions. \(\iota\) requires no prior cause: it is the minimum condition for anything to be conceivable at all.

A single \(\iota\) does not carry character or polarity. It produces a split: two complementary aspects \(\{A,\, \bar{A}\}\) that are not each other. Neither \(A\) nor \(\bar{A}\) is intrinsically positive or negative. Polarity is not a property of a single \(\iota\).

Let \(D\) denote the collection of all \(\iota\) distinctions within \(M\).

\subsection{Character from Comparison}
\label{sec:org0a09699}

When two distinctions \(i, j \in D\) co-exist within \(M\), a comparison is unavoidable: do they distinguish in the same way, or in opposite ways? This comparison is binary and symmetric:

\begin{equation}
c : D \times D \to \{\mathrm{same},\, \mathrm{different}\}
\end{equation}

\(c(i,i) = \mathrm{same}\). \(c(i,j) = c(j,i)\). What earlier formulations labeled \(+1\) or \(-1\) is the outcome of this comparison. Character and polarity are the first emergent properties of the substrate, not primitives.

\subsection{The Relational Web}
\label{sec:org8469765}

Co-existing distinctions in \(D\) necessarily relate. When \(D\) contains more than one distinction, comparisons arise throughout, generating a relational graph \(G\):

\begin{itemize}
\item Nodes: \(\iota\) distinctions in \(D\)
\item Edges: relations between co-existing distinctions
\end{itemize}

The negotiation field \(N_{ij}\) measures the coherence of the comparison between distinctions \(i\) and \(j\) across the entire relational structure --- not the comparison value itself, but the degree to which that comparison is consistent across all paths through \(G\). \(N\) is real-valued and continuous: \(N_{ij} > 0\) (different-comparing, stable); \(N_{ij} < 0\) (same-comparing, unstable); \(N_{ij} = 0\) (uncorrelated).

The Laplacian \(L\) of \(G\) encodes the relational geometry. In the large-scale continuum limit \(L\) converges to the Laplace-Beltrami operator on a Riemannian manifold --- geometry emerging from the relational structure of distinctions, not assumed.

\subsection{The Adjacency Constraint}
\label{sec:org3b92a93}

Two distinctions comparing as \(\mathrm{same}\) carry no new relational information: their comparison is trivially resolved. Negotiations between same-comparing distinctions contribute nothing to stable structure. This is not a topological prohibition on edges --- it is a statement about information content. Productive negotiation requires \(c(i,j) = \mathrm{different}\). Stable configurations require all local comparisons to be productive.

\subsection{A Note on Time}
\label{sec:org25e0df2}

The dynamics below use \(\partial N / \partial t\). Time appears here as a background parameter retained for calculability. In the substrate as described, sequence does not exist prior to distinction --- time is not a primitive but an emergent ordering of the relational structure of \(D\). The derivation of temporal sequence from substrate dynamics is an open problem.

\section{Scale Inputs}
\label{sec:orgbb3add4}

This description has exactly two inputs:

\begin{center}
\begin{tabular}{lll}
\toprule
\textbf{Input} & \textbf{Value} & \textbf{Role}\\[0pt]
\midrule
Iota scale & \(l_\iota = 1.616 \times 10^{-35}\) m & Substrate grain\\[0pt]
Hubble scale & \(R_H = c/H_0\) (current epoch) & Longest negotiation mode\\[0pt]
\bottomrule
\end{tabular}
\end{center}

In Planck units: \(G = 1\), \(c = 1\), \(\hbar = 1/(2\pi)\). What SI physics calls \(\hbar\) is the measurement of one projected \(\iota\) event in the observer's unit system --- a conversion artifact, not a primitive constant.

The negotiation energy quantum is estimated from the geometric mean of the two scale energies:

\begin{equation}
\varepsilon_\text{est} = \bigl(E_\text{Planck}^2 \cdot E_\text{Hubble}\bigr)^{1/3} \approx 152.6\ \text{MeV}
\end{equation}

With a \((2\pi)^{1/3}\) dispersion factor this reaches 281.5 MeV. The observed value \(\varepsilon = m_\text{proton}/3 = 312.8\) MeV is a definition from observation; the \(\sim\)10$\backslash$% residual gap is closed by the Barbero-Immirzi correction (\(\gamma \approx 0.2375\)), though that correction has not yet been derived from the negotiation graph spectrum.

\section{Dynamics}
\label{sec:org652e8c2}

The negotiation field \(N_{ij}\) encodes the relational coherence between adjacent \(\iota\) distinctions \(i\) and \(j\). Lagrangian in Planck units:

\begin{equation}
\mathcal{L} = \tfrac{1}{2}\!\left(\frac{\partial N}{\partial t}\right)^{\!2} - \tfrac{1}{2}\,N^\top L\, N
\end{equation}

Hamiltonian (via Legendre transform, \(\Pi = \partial N/\partial t\)):

\begin{equation}
\hat{H} = \tfrac{1}{2}\Pi^2 + \tfrac{1}{2}\,N^\top L\, N
\end{equation}

Energy eigenvalues: \(E_k = \varepsilon\,\lambda_k\), where \(\lambda_k\) are eigenvalues of the graph Laplacian \(L\).

Stable particles satisfy the closure condition \(L \cdot N = 0\). For a complete graph on \(n\) nodes with mode number \(k\), energy eigenvalues follow:

\begin{equation}
E(n,k) = \varepsilon \cdot k \cdot n
\end{equation}

The minimum closed graph is the triangle of three \(\iota\) distinctions (\(\lambda = 3\), \(n = 3\)):

\begin{equation}
m_\text{proton} = 3\varepsilon = 938.3\ \text{MeV}
\end{equation}

Higher topologies follow the same formula: a two-node graph corresponds to the rho meson (\(n=2\), predicted 625.6 MeV); a four-node graph corresponds to the tetraquark (\(n=4\), predicted 1251 MeV). The triangle is the ground state --- the minimum topology that achieves closure.

\section{Derived Results}
\label{sec:orgb6cd773}

From two inputs and the single correction factor \(\gamma \approx 0.2375\):

\begin{center}
\begin{tabular}{llll}
\toprule
\textbf{Quantity} & \textbf{Derived} & \textbf{Measured} & \textbf{Agreement}\\[0pt]
\midrule
\(m_\text{proton}\) & \(3\varepsilon = 938.3\) MeV & 938.3 MeV & Exact (definition of \(\varepsilon\))\\[0pt]
\(G\) & \(H_0^2 / 4\pi\rho_\Lambda\) & \(6.674\times10^{-11}\) SI & 94--97\%\\[0pt]
\(\hbar\) & \(m_p c\,k(\gamma)\,r_p / 6\pi\) & \(1.055\times10^{-34}\) J\(\cdot\)s & 99.1\%\\[0pt]
\(H_0\) & 68.6 km/s/Mpc & 67.4--73.0 (tension range) & Within tension\\[0pt]
\(\Lambda\) & \(\approx 10^{-122}\) Planck & \(\approx 10^{-122}\) Planck & Same order\\[0pt]
\bottomrule
\end{tabular}
\end{center}

The ratio \(\rho_\text{Planck}/\rho_\Lambda \approx 10^{123}\) is the depth of the Medium --- the dynamic range from iota grain to cosmic horizon --- not a fine-tuning problem.

Analytic prediction: \(\Omega_\Lambda = 2/3 \approx 0.667\) (observed: 0.685; discrepancy 2.7$\backslash$%, attributable to matter contribution). Testable from expansion history without measuring \(G\).

All three residual gaps (in \(G\), \(m_\text{proton}\), and \(\hbar\)) are closed by a single Barbero-Immirzi calculation on the negotiation graph. This is the primary open calculation in this description.

\section{License and Attribution}
\label{sec:org574c6d4}

CC BY 4.0. DOI: 10.5281/zenodo.18927154 \quad ORCID: 0009-0001-0328-2299 \quad github.com/bkdepuy/The-Medium

\section{Document Series and Cross-References}
\label{sec:org08041df}

This is Document 1 of 5 in The Medium series.

\begin{center}
\begin{tabular}{clp{9cm}}
\toprule
\textbf{Doc} & \textbf{File} & \textbf{Content}\\[0pt]
\midrule
0 & medium\_doc0\_introduction.org & Personal introduction; the question behind the description\\[0pt]
1 & medium\_doc1\_definition.org & \emph{This paper}\\[0pt]
2 & medium\_doc2\_unified.org & Negotiation Modes, Particle Structure, Proton Mass\\[0pt]
3 & medium\_doc3\_lagrangian.org & Lagrangian of the Negotiation Field; QFT connection; SU(3)\\[0pt]
4 & medium\_doc4\_hamiltonian.org & Hamiltonian, energy spectrum, canonical quantization, ADM\\[0pt]
5 & medium\_doc5\_G\_H0.org & Deriving \(G\) and \(H_0\) from substrate primitives\\[0pt]
\bottomrule
\end{tabular}
\end{center}
\end{document}
